% THEOREM 5.14 ARCHIVE: Verbatim theorem-and-proof source from the rejected manuscript; candidate for a separate note and intentionally not input by main.tex.
\subsection{Optimal injective relabeling of the realized Sturmian law}

\begin{theorem}[Optimal injective relabeling of the realized Sturmian law]
\label{thm:optimal-injective-relabeling}
Assume \(\beta=\alpha\). For \(m\ge1\), let
\[
\mathcal I_m(\alpha):=
\{\iota:S_m(\alpha)\to X_m:\ \iota\text{ is injective}\},
\]
and for \(\iota\in\mathcal I_m(\alpha)\) write
\[
\nu_\iota:=\iota_*\mu_m.
\]
If \(m\ge6\), then for every \(q\in[0,\infty]\),
\[
\inf_{\iota\in\mathcal I_m(\alpha)}D_q(\nu_\iota\|\pi_m)
=
(m-1)\log\varphi-\log\pi(0)-H_q(\mu_m).
\]
Every injection with image contained in \(X_m^{00}\) attains this value. For
every \(q<\infty\), these are exactly the optimal injections. For
\(q=\infty\), write
\[
p_*:=\max_{w\in S_m(\alpha)}\mu_m(w),\qquad
\tau(u):=\frac{\pi_m(u)}{\pi(0)\varphi^{-(m-1)}}\in\{1,\varphi^{-1},\varphi^{-2}\}.
\]
An injection \(\iota\) is optimal if and only if
\[
\mu_m(w)\le p_*\tau(\iota(w))
\qquad\text{for every }w\in S_m(\alpha).
\]
Thus \(q=\infty\) may have additional optimizers, but precisely the additional
placements of non-maximal source atoms whose loss of Parry mass is no larger
than their drop below the maximal source mass. Since \(|X_m^{00}|=F_m\), an
all-\(X_m^{00}\) optimal placement exists if and only if \(m\ge6\).
\end{theorem}

\begin{proof}
Fix \(\iota\in\mathcal I_m(\alpha)\), and put \(\nu:=\nu_\iota\). Since
\(\iota\) is injective on the \(m+1\)-point Sturmian support \(S_m(\alpha)\),
it only relabels the atoms of \(\mu_m\). Therefore
\[
H_q(\nu)=H_q(\mu_m)
\qquad (q\in[0,\infty]).
\]

Set
\[
c_{00}:=\pi(0)\varphi^{-(m-1)}.
\]
By Proposition~\ref{prop:parry-endpoint-collapse},
\[
\pi_m(u)=c_{00}\tau(u),
\qquad
\tau(u)\in\{1,\varphi^{-1},\varphi^{-2}\},
\]
and \(\tau(u)=1\) exactly on \(X_m^{00}\).

First let \(q\in(0,\infty)\setminus\{1\}\). Then
\[
\begin{aligned}
D_q(\nu\|\pi_m)
&=
\frac{1}{q-1}
\log\!\left(\sum_{u\in X_m}\nu(u)^q(c_{00}\tau(u))^{1-q}\right)\\
&=
-\log c_{00}-H_q(\nu)
+
\frac{1}{q-1}
\log\!\left(
\frac{\sum_{u\in X_m}\nu(u)^q\tau(u)^{1-q}}
{\sum_{u\in X_m}\nu(u)^q}
\right).
\end{aligned}
\]
Because \(0<\tau(u)\le1\) for every \(u\), the last term is nonnegative. It
vanishes if and only if \(\tau(u)=1\) on \(\supp(\nu)\), equivalently
\(\supp(\nu)\subset X_m^{00}\). Hence
\[
D_q(\nu\|\pi_m)
\ge
(m-1)\log\varphi-\log\pi(0)-H_q(\mu_m),
\]
with equality precisely when \(\iota(S_m(\alpha))\subset X_m^{00}\).

For \(q=1\), Theorem~\ref{thm:parry-kl-identity} gives
\[
\begin{aligned}
\KL(\nu\|\pi_m)
&=
m\log\varphi-\log\pi(0)-H_1(\nu)
+
\bigl(\eta_{11}(\nu)-\eta_{00}(\nu)\bigr)\log\varphi\\
&=
(m-1)\log\varphi-\log\pi(0)-H_1(\mu_m)\\
&\quad
+
\bigl(\eta_{01}(\nu)+\eta_{10}(\nu)+2\eta_{11}(\nu)\bigr)\log\varphi.
\end{aligned}
\]
The correction term is nonnegative and vanishes exactly when
\(\eta_{01}(\nu)=\eta_{10}(\nu)=\eta_{11}(\nu)=0\), that is, exactly when
\(\supp(\nu)\subset X_m^{00}\).

For \(q=0\), \(\nu\) has support size \(m+1\), so
\[
D_0(\nu\|\pi_m)
=
-\log\pi_m(\supp\nu)
\ge
-\log\bigl((m+1)c_{00}\bigr)
=
(m-1)\log\varphi-\log\pi(0)-H_0(\mu_m).
\]
Equality holds if and only if every point of \(\supp(\nu)\) has Parry mass
\(c_{00}\), namely when \(\supp(\nu)\subset X_m^{00}\).

For \(q=\infty\),
\[
\begin{aligned}
D_\infty(\nu\|\pi_m)
&=
\log\max_{u\in\supp(\nu)}\frac{\nu(u)}{\pi_m(u)}\\
&\ge
\log\max_{u\in\supp(\nu)}\frac{\nu(u)}{c_{00}}
=
-\log c_{00}-H_\infty(\nu)\\
&=
(m-1)\log\varphi-\log\pi(0)-H_\infty(\mu_m).
\end{aligned}
\]
The displayed lower bound is sharp exactly when
\[
\max_{u\in\supp(\nu)}\frac{\nu(u)}{c_{00}\tau(u)}
=
\frac{\max_{u\in\supp(\nu)}\nu(u)}{c_{00}}.
\]
Since \(\nu=\iota_*\mu_m\) and \(\iota\) is injective, this condition is
equivalent to
\[
\frac{\mu_m(w)}{\tau(\iota(w))}\le p_*
\qquad (w\in S_m(\alpha)),
\]
or, equivalently, to
\(\mu_m(w)\le p_*\tau(\iota(w))\) for every realized word \(w\).  This proves
the asserted \(q=\infty\) optimizer characterization.  In particular, any
injection with image in \(X_m^{00}\), for which \(\tau\equiv1\) on the image,
attains equality here as well; other optimizers occur exactly when every atom
placed outside \(X_m^{00}\) is small enough to satisfy the displayed
inequality.

It remains to characterize the admissible range. The Sturmian complexity is
\(|S_m(\alpha)|=m+1\), while Lemma~\ref{lem:endpoint-class-counts} gives
\(|X_m^{00}|=F_m\). Since \(F_5=5<6\) and \(F_6=8\ge7\), and
\[
F_{m+1}-(m+2)=\bigl(F_m-(m+1)\bigr)+F_{m-1}-1,
\]
an induction on \(m\) shows that \(F_m\ge m+1\) exactly for \(m\ge6\). Thus
an injective placement inside \(X_m^{00}\) exists exactly in that range, and
the displayed infimum formula follows.
\end{proof}
